
\subsection{Dual-Path Networks}
\label{sec:related-dpn}
The original DPRNN \cite{dprnn} architecture was proposed for waveform-domain speech separation. It divides an input sequence of $C$ channels and $L$ samples into short chunks with a chunk size of $K$ and a hop size of $P$, resulting in a three-dimensional tensor of shape $[C, L/P, K]$. The model then iteratively performs an intra-chunk RNN along the third dimension and an inter-chunk RNN along the second dimension. This effectively models long sequences and outperforms previous waveform-domain approaches. DPRNN has also been widely adopted for time-frequency (TF)-domain speech enhancement \cite{dpcrn, trunet, fspen, gtcrn, lisennet, bsrnn-se}. It processes a spectrogram by alternating between a time RNN that models temporal dependencies and a frequency RNN that models dependencies across frequency bands. However, applying an RNN to the frequency axis assumes that frequency bands have a sequential relationship. This is suboptimal because spectral relationships (e.g. harmonics) exhibit global rather than sequential dependencies. We argue that a model should be capable of handling all frequency interactions concurrently.

Some prior works have attempted to replace the RNNs in dual-path networks with other powerful modules like Mamba and xLSTM \cite{semamba, xlstm-senet}. However, these models are designed for non-streaming, large-scale scenarios, and still impose a sequential relationship on the frequency axis. Meanwhile, others have explored dual-path transformers (DPT) \cite{dptnet, sepformer}. Transformers are intuitively well-suited for modeling global information across frequency bands. On the other hand, a time-axis transformer is problematic for low-latency streaming. It requires caching previous key and value frames, demanding additional memory operations. Moreover, to improve speech quality, the number of cache frames should be increased \cite{streaming-dpt}, requiring more memory and arithmetic operations.

In this paper, we propose RNNFormer, which is composed of a time RNN and a frequency transformer (Section \ref{sec:model}). The time RNN maintains the low-complexity streaming nature of our model, while the frequency transformer captures global information across all frequency bands. We experimentally prove that RNNFormers achieve a superior RTF-performance curve compared to DPRNNs and DPTs (Section \ref{sec:ablation}).

\subsection{Processing Speed of Speech Enhancement Models}
\label{sec:related-speed}
Many prior works have focused on reducing the number of parameters and MACs to improve model efficiency. DPCRN \cite{dpcrn} and TRU-Net \cite{trunet} introduced a TF-domain U-Net structure with a convolutional encoder, decoder, and DPRNNs. Building upon this, more elaborate methods such as sub-band processing and enhanced DPRNN structures have emerged.

BSRNN \cite{bsrnn-music-separation} introduced a band-split encoder and decoder. FSPEN \cite{fspen} proposed a sub-band encoder and decoder and a time RNN with a path extension. GTCRN \cite{gtcrn} and LiSenNet \cite{lisennet} utilized sub-band feature extraction, grouped DPRNNs, and channel split-concatenate-shuffle operations. While these methods successfully reduced the number of parameters and MACs, their increased architectural complexity and memory operations paradoxically led to higher RTFs (Section \ref{sec:performance}).

In contrast to these approaches, we directly aim to minimize the RTF for practical scenarios. We experimentally prove that a simple network design -- without complex sub-band operations or RNN grouping -- is sufficient to obtain a powerful and fast model.